% titel: Blatt 6, A4: Überlagerte Schwingungen
% Teil II: der durchgerechnete Weg.

\subsection*{Gegeben und gesucht}
$A>0$, $\varphi\in[0,2\pi)$ finden mit
\begin{equation*}
  \cos(\omega t)+2\cos\!\Bigl(\omega t+\frac{2\pi}{3}\Bigr)
  =A\cos(\omega t+\varphi).
\end{equation*}

\subsection*{Ansatz und Gedankenschritt}
\emph{Warum Zeiger?} Beide Schwingungen haben dieselbe Kreisfrequenz
$\omega$; der gemeinsame Faktor $\mathrm e^{i\omega t}$ ist überall
gleich. Es bleibt, die beiden (fixen) komplexen Amplituden zu addieren und
das Ergebnis in Polarkoordinaten zu lesen -- kein mühsames
Trigonometrie-Umrechnen nötig.

\subsection*{Zeigersumme}
Die Schwingungen als komplexe Amplituden:
\begin{equation*}
  z_1=\mathrm e^{i0}=1,\qquad
  z_2=2\,\mathrm e^{i\frac{2\pi}{3}}
  =2\Bigl(-\frac12+i\frac{\sqrt3}{2}\Bigr)=-1+i\sqrt3 .
\end{equation*}
\emph{Gedankenschritt:} Der Zeiger $z_2$ zeigt vom Ursprung in Richtung
$120^\circ$ mit Länge $2$; in Komponenten heisst das $(-1,+\sqrt3)$. Der
erste Zeiger ist der Einheitszeiger längs der positiven $x$-Achse, die
sich mit $z_2$ zu \emph{genau} $i\sqrt3$ addieren (die $x$-Komponenten
heben sich auf).
\begin{equation*}
  z=z_1+z_2=1+(-1+i\sqrt3)=i\sqrt3 .
\end{equation*}
In Polarkoordinaten: $A=|z|=\sqrt3$ und $\varphi=\arg(z)=\frac\pi2$.

\subsection*{Rückübersetzung}
Die Zeigersumme $i\sqrt3$ zurück in einen Kosinus:
\begin{equation*}
  \cos(\omega t)+2\cos\!\Bigl(\omega t+\frac{2\pi}{3}\Bigr)
  =\sqrt3\,\cos\!\Bigl(\omega t+\frac\pi2\Bigr).
\end{equation*}

\begin{ergebnis}
\begin{equation*}
  A=\sqrt3,\qquad \varphi=\frac\pi2 .
\end{equation*}
\end{ergebnis}

\begin{kontrolle}
\textbf{Probe (Zeitachse):} Für $\omega t=-\frac\pi2$ ist die rechte Seite
$\sqrt3\cos(0)=\sqrt3$; links:
$\cos(-\frac\pi2)+2\cos(\frac\pi6)=0+2\cdot\frac{\sqrt3}{2}=\sqrt3$.
\checkmark Für $\omega t=0$: links $1+2\cos\frac{2\pi}{3}=1-1=0$, rechts
$\sqrt3\cos\frac\pi2=0$. \checkmark

\textbf{Grenzfall:} Wären beide Zeiger in dieselbe Richtung gerichtet
(gleiche Phase), ergäbe sich $A=1+2=3$ (konstruktive Interferenz); hier
heben sich die $x$-Anteile weg (Phasendifferenz wirkt), daher die kleinere
Amplitude $\sqrt3$ bei $90^\circ$ Phasenverschiebung. \checkmark
\end{kontrolle}

\begin{merksatz}
Gleichfrequente Kosinus-Schwingungen addieren sich über ihre komplexen
Amplituden (Zeiger): Summe der Zeiger $\sum a_j\mathrm e^{i\varphi_j}$
bilden, dann $A$ (Betrag der Zeigersumme), $\varphi$ (deren Argument) ablesen. Das Ergebnis ist
wieder ein Kosinus \emph{derselben} Frequenz. Verschiedene Phasen
verstärken oder abschwächen die resultierende Amplitude.
\end{merksatz}
